\documentclass[10pt,a4paper]{amsart}
\usepackage{amsmath,amssymb,amsthm,graphicx}
% Retained minimal preamble for the Stacks Project chapter more-morphisms.tex (More on Morphisms).
% Verbatim excerpt of the chapter preamble: Xy-pic setup (native preamble lines 14-18) and the native
% theorem declarations with the mathematical macros (native preamble lines 154-217). Copied mechanically;
% nothing added, removed or redefined. The wrapper class and the amsmath/amssymb/amsthm/graphicx base
% packages are supplied by the assembler, so the class/externaldocument part of the native preamble is omitted.
% For commutative diagrams we use Xy-pic
\usepackage[all]{xy}

% We use 2cell for 2-commutative diagrams.
\xyoption{2cell}
% Theorem environments.
%
\theoremstyle{plain}
\newtheorem{theorem}[subsection]{Theorem}
\newtheorem{proposition}[subsection]{Proposition}
\newtheorem{lemma}[subsection]{Lemma}

\theoremstyle{definition}
\newtheorem{definition}[subsection]{Definition}
\newtheorem{example}[subsection]{Example}
\newtheorem{exercise}[subsection]{Exercise}
\newtheorem{situation}[subsection]{Situation}

\theoremstyle{remark}
\newtheorem{remark}[subsection]{Remark}
\newtheorem{remarks}[subsection]{Remarks}

\numberwithin{equation}{subsection}

% Macros
%
\def\lim{\mathop{\mathrm{lim}}\nolimits}
\def\colim{\mathop{\mathrm{colim}}\nolimits}
\def\Spec{\mathop{\mathrm{Spec}}}
\def\Hom{\mathop{\mathrm{Hom}}\nolimits}
\def\Ext{\mathop{\mathrm{Ext}}\nolimits}
\def\SheafHom{\mathop{\mathcal{H}\!\mathit{om}}\nolimits}
\def\SheafExt{\mathop{\mathcal{E}\!\mathit{xt}}\nolimits}
\def\Sch{\mathit{Sch}}
\def\Mor{\mathop{\mathrm{Mor}}\nolimits}
\def\Ob{\mathop{\mathrm{Ob}}\nolimits}
\def\Sh{\mathop{\mathit{Sh}}\nolimits}
\def\NL{\mathop{N\!L}\nolimits}
\def\CH{\mathop{\mathrm{CH}}\nolimits}
\def\proetale{{pro\text{-}\acute{e}tale}}
\def\etale{{\acute{e}tale}}
\def\QCoh{\mathit{QCoh}}
\def\Ker{\mathop{\mathrm{Ker}}}
\def\Im{\mathop{\mathrm{Im}}}
\def\Coker{\mathop{\mathrm{Coker}}}
\def\Coim{\mathop{\mathrm{Coim}}}

% Boxtimes
%
\DeclareMathSymbol{\boxtimes}{\mathbin}{AMSa}{"02}

%
% Macros for moduli stacks/spaces
%
\def\QCohstack{\mathcal{QC}\!\mathit{oh}}
\def\Cohstack{\mathcal{C}\!\mathit{oh}}
\def\Spacesstack{\mathcal{S}\!\mathit{paces}}
\def\Quotfunctor{\mathrm{Quot}}
\def\Hilbfunctor{\mathrm{Hilb}}
\def\Curvesstack{\mathcal{C}\!\mathit{urves}}
\def\Polarizedstack{\mathcal{P}\!\mathit{olarized}}
\def\Complexesstack{\mathcal{C}\!\mathit{omplexes}}
% \Pic is the operator that assigns to X its picard group, usage \Pic(X)
% \Picardstack_{X/B} denotes the Picard stack of X over B
% \Picardfunctor_{X/B} denotes the Picard functor of X over B
\def\Pic{\mathop{\mathrm{Pic}}\nolimits}
\def\Picardstack{\mathcal{P}\!\mathit{ic}}
\def\Picardfunctor{\mathrm{Pic}}
\def\Deformationcategory{\mathcal{D}\!\mathit{ef}}
\usepackage{cleveref}
\crefname{originallemma}{original lemma}{original lemmas}
\Crefname{originallemma}{Original lemma}{Original lemmas}
\makeatletter
\newlabel{lemma-pushout-functor}{{0ECK}{source}{}{sourcepointer-1}{}}
\newlabel{lemma-pushout-functor@cref}{{[originallemma][0][]0ECK}{[1][0][]source}{}{}{}}
\makeatother
\crefname{originalreference}{original reference}{original references}
\Crefname{originalreference}{Original reference}{Original references}
\makeatletter
\newlabel{more-algebra-lemma-module-over-fibre-product-bis}{{08IG}{source}{}{sourcepointer-2}{}}
\newlabel{more-algebra-lemma-module-over-fibre-product-bis@cref}{{[originalreference][0][]08IG}{[1][0][]source}{}{}{}}
\makeatother
\crefname{originalreference}{original reference}{original references}
\Crefname{originalreference}{Original reference}{Original references}
\makeatletter
\newlabel{more-algebra-lemma-properties-algebras-over-fibre-product}{{08KQ}{source}{}{sourcepointer-3}{}}
\newlabel{more-algebra-lemma-properties-algebras-over-fibre-product@cref}{{[originalreference][0][]08KQ}{[1][0][]source}{}{}{}}
\makeatother
\crefname{originalreference}{original reference}{original references}
\Crefname{originalreference}{Original reference}{Original references}
\makeatletter
\newlabel{situation-pushout-along-closed-immersion-and-integral}{{0ECI}{source}{}{sourcepointer-4}{}}
\newlabel{situation-pushout-along-closed-immersion-and-integral@cref}{{[originalreference][0][]0ECI}{[1][0][]source}{}{}{}}
\makeatother
\title{Flat and etale descent along the pushout square: the affine flip equivalence}
\author{The Stacks Project Authors}
\date{stacks-project:more-morphisms.tex:local-snapshot-20261011; GNU FDL 1.2 or later}
\begin{document}
\maketitle
\noindent\textit{Editorial note.} Counted span is the author's own proof environment, source lines 19955-19980, immediately following the statement at 19945-19953.\par\medskip
\noindent\textit{Editorial note.} The sentence recording that the comparison arrow need not be an isomorphism without flatness (source lines 19965-19971) is part of the counted proof and is the reason the flat hypothesis is essential.\par\medskip
\noindent\textit{Editorial note.} Two cross-chapter pointer families (More on Algebra) are recorded in assembly.external\_references; the Situation and the defining Lemma are same-chapter line-indexed stubs.\par\medskip
\noindent\textit{Allowed standard prerequisites (editorial).}\begin{itemize}
\item The pushout situation and its functor. Situation situation-pushout-along-closed-immersion-and-integral and Lemma lemma-pushout-functor of this chapter, cited at statement lines 19945-19948. Fixes the pushout diagram Y amalg\_Z X (Z a closed subscheme of X with X -> Y integral) and the data tuple (X', Y', Z', i', j', f, g, h) whose category is being compared.
\item Flatness and etaleness of the induced pushout morphism. More on Algebra, Lemma more-algebra-lemma-properties-algebras-over-fibre-product, cited at proof lines 19952-19954 and 19975-19977. Both directions hinge on it: forward, Y' amalg\_\{Z'\} X' -> Y amalg\_Z X is flat/etale; converse, the comparison arrow is an isomorphism exactly when W is flat.
\item Closed immersion behaviour without flatness. More on Algebra, Lemma more-algebra-lemma-module-over-fibre-product-bis, cited at proof lines 19968-19971. Records that without flatness the comparison arrow is only a closed immersion in general, which is why the flat hypothesis cannot be dropped.
\end{itemize}\medskip
\ifdefined\hypertarget\hypertarget{sourcepointer-1}{}\fi
\noindent\textit{Source reference.} Original lemma 0ECK: sources/stacks-more-morphisms/pinned-original-source.tex, source lines 19876--19876. Navigation pointer only: the counted proof cites this statement by \textbackslash{}ref as a standing input; the deduction the counted proof has to make is carried out inside the counted span, and no part of that deduction is delegated to this pointer.\par\medskip
\ifdefined\hypertarget\hypertarget{sourcepointer-2}{}\fi
\noindent\textit{Source reference.} Original reference 08IG: focused-seeds/geometry/S9-morph-pushout-functor-equivalence-flat/ref-tags.txt, source lines 5--5. Navigation pointer only: the counted proof cites this statement by \textbackslash{}ref as a standing input; the deduction the counted proof has to make is carried out inside the counted span, and no part of that deduction is delegated to this pointer.\par\medskip
\ifdefined\hypertarget\hypertarget{sourcepointer-3}{}\fi
\noindent\textit{Source reference.} Original reference 08KQ: focused-seeds/geometry/S9-morph-pushout-functor-equivalence-flat/ref-tags.txt, source lines 6--6. Navigation pointer only: the counted proof cites this statement by \textbackslash{}ref as a standing input; the deduction the counted proof has to make is carried out inside the counted span, and no part of that deduction is delegated to this pointer.\par\medskip
\ifdefined\hypertarget\hypertarget{sourcepointer-4}{}\fi
\noindent\textit{Source reference.} Original reference 0ECI: sources/stacks-more-morphisms/pinned-original-source.tex, source lines 19708--19708. Navigation pointer only: the counted proof cites this statement by \textbackslash{}ref as a standing input; the deduction the counted proof has to make is carried out inside the counted span, and no part of that deduction is delegated to this pointer.\par\medskip
\begin{lemma}
\label{lemma-pushout-functor-equivalence-flat}
In Situation \ref{situation-pushout-along-closed-immersion-and-integral}
the category of schemes flat, separated, and locally quasi-finite
over the pushout $Y \amalg_Z X$ is equivalent to the category of
$(X', Y', Z', i', j', f, g, h)$ as in Lemma \ref{lemma-pushout-functor}
with $f, g, h$ flat. Similarly with ``flat'' replaced with
``\'etale''.
\end{lemma}

\begin{proof}
If we start with $(X', Y', Z', i', j', f, g, h)$
as in Lemma \ref{lemma-pushout-functor} with $f, g, h$ flat
or \'etale, then $Y' \amalg_{Z'} X' \to Y \amalg_Z X$ is flat
or \'etale by More on Algebra, Lemma
\ref{more-algebra-lemma-properties-algebras-over-fibre-product}.

\medskip\noindent
For the converse, let
$W \to Y \amalg_Z X$ be a separated and locally quasi-finite morphism.
Set $X' = W \times_{Y \amalg_Z X} X$, $Y' = W \times_{Y \amalg_Z X} Y$, and
$Z' = W \times_{Y \amalg_Z X} Z$ with obvious morphisms $i', j', f, g, h$.
Form the pushout $Y' \amalg_{Z'} X'$. We obtain a morphism
$$
Y' \amalg_{Z'} X' \longrightarrow W
$$
of schemes over $Y \amalg_X Z$ by the universal property of the
pushout. If we do not assume that $W \to Y \amalg_Z X$ is flat,
then in general this morphism won't be an isomorphism.
(In fact, More on Algebra, Lemma
\ref{more-algebra-lemma-module-over-fibre-product-bis}
shows the displayed arrow is a closed immersion but not an isomorphism
in general.) However, if $W \to Y \times_Z X$ is flat, then
it is an isomorphism by More on Algebra, Lemma
\ref{more-algebra-lemma-properties-algebras-over-fibre-product}.
\end{proof}

\end{document}
